\chapter{Menyelesaikan Soal Cerita}\label{ch:g2:problems}

Sebuah cerita yang berisi bilangan menyembunyikan satu pertanyaan:
menjumlahkan? mengurangkan? mengalikan? Bab ini memberi cara untuk
memutuskan, memakai gambar batang sederhana agar operasinya
\emph{terlihat}, lalu membaca keterangan dari tabel dan piktogram.

\section{Memilih operasinya}

\begin{method}[Operasi yang mana?]\label{met:g2:problems:choose}
\begin{enumerate}
  \item \emph{digabungkan, seluruhnya}: jumlahkan ($+$);
  \item \emph{diambil, tinggal berapa, berapa lebihnya}: kurangkan
        ($-$);
  \item \emph{sekian kelompok yang sama banyak}: kalikan ($\times$);
  \item gambarlah sketsa cepat kalau ragu, dan jawablah selalu
        dengan satu kalimat.
\end{enumerate}
\end{method}

\begin{omfigure}
\begin{tikzpicture}[scale=0.6]
  % digabungkan
  \draw[thick, fill=omDef!20] (0,0) rectangle (3,0.8);
  \draw[thick, fill=omThm!20] (3,0) rectangle (4.6,0.8);
  \node[font=\small] at (1.5,0.4) {$24$};
  \node[font=\small] at (3.8,0.4) {$15$};
  \draw[Stealth-Stealth] (0,-0.4) -- (4.6,-0.4);
  \node[below, font=\small] at (2.3,-0.5) {? seluruhnya: $24 + 15$};
  % dibandingkan
  \begin{scope}[xshift=8cm]
  \draw[thick, fill=omDef!20] (0,0.9) rectangle (4.6,1.7);
  \node[font=\small] at (2.3,1.3) {$40$};
  \draw[thick, fill=omThm!20] (0,0) rectangle (2.8,0.8);
  \node[font=\small] at (1.4,0.4) {$28$};
  \draw[Stealth-Stealth] (2.8,0.4) -- (4.6,0.4);
  \node[below, font=\small] at (2.3,-0.2) {? lebihnya: $40 - 28$};
  \end{scope}
\end{tikzpicture}

{\small Gambar batang: dua batang yang disambung berarti penjumlahan;
dua batang yang dibandingkan berarti pengurangan.}
\end{omfigure}

\begin{example}[Soal dua langkah]\label{ex:g2:problems:twostep}
``Sandi membeli $3$ bungkus berisi $6$ stiker, lalu memberikan $5$
kepada temannya. Berapa yang ia simpan?''
\begin{enumerate}
  \item Kelompok: $3 \times 6 = 18$ stiker;
  \item diberikan: $18 - 5 = 13$.
\end{enumerate}
Sandi menyimpan $13$ stiker. Bacalah lagi pertanyaannya: apakah
ceritanya benar-benar sudah selesai? Di sini sudah.
\end{example}

\begin{example}[Bilangan tambahan yang tersembunyi]\label{ex:g2:problems:extra}
``Sebuah buku setebal $210$ halaman\,\dots'' --- hati-hati, tidak
setiap bilangan dalam cerita itu diperlukan. ``Beni punya $8$
kelereng merah dan $6$ kelereng biru, dan umurnya $7$ tahun. Berapa
kelerengnya?'' Umur itu umpan: $8 + 6 = 14$ kelereng. Pilahlah
bilangannya sebelum menghitung.
\end{example}

\section{Tabel dan piktogram}

\begin{example}[Sebuah tabel]\label{ex:g2:problems:table}
Kue yang terjual di sebuah kios:
\begin{center}
\renewcommand{\arraystretch}{1.3}
\begin{tabular}{l|ccc}
 & Sabtu & Minggu & \omterm{def:g1:addition:def}{jumlah} \\
\hline
muffin & $12$ & $9$ & $21$ \\
pai & $7$ & $10$ & $17$ \\
\end{tabular}
\end{center}
Membaca: pai pada hari Minggu, $10$. Menghitung: muffin selama dua
hari, $12 + 9 = 21$; pai pada hari Minggu lebih banyak daripada hari Sabtu sebanyak $10 - 7 = 3$.
\end{example}

\begin{example}[Sebuah piktogram]\label{ex:g2:problems:pictogram}
Bacalah kuncinya dulu: di sini satu apel mewakili $2$ buah.

\begin{omfigure}
\begin{tikzpicture}[scale=0.62]
  \node[left, font=\small] at (0,2) {merah};
  \foreach \i in {0,1,2,3} \fill[omThm] (0.5+\i*0.75,2) circle (0.24);
  \node[left, font=\small] at (0,1) {hijau};
  \foreach \i in {0,1} \fill[omMeth] (0.5+\i*0.75,1) circle (0.24);
  \node[left, font=\small] at (0,0) {kuning};
  \foreach \i in {0,1,2} \fill[omProp] (0.5+\i*0.75,0) circle (0.24);
  \node[right, font=\small] at (4.0,1) {kunci: satu titik $= 2$ buah};
\end{tikzpicture}

{\small Sebuah piktogram: merah menunjukkan $4$ titik, jadi
$4 \times 2 = 8$ buah merah; hijau $2$ titik, $4$ buah; kuning $3$
titik, $6$ buah.}
\end{omfigure}
\end{example}

\section{Latihan}

\begin{exercise}[$\star$]\label{exo:g2:problems:1}
Sebutkan operasinya (jangan dihitung): ``$34$ dan $28$ kelereng
digabungkan''; ``$50$ permen, $12$ dimakan''; ``$4$ kotak berisi $5$
apel''; ``Ani punya $30$, Leo punya $18$: berapa lebihnya Ani?''.
\end{exercise}

\begin{exercise}[$\star$]\label{exo:g2:problems:2}
Selesaikan: ``Sebuah tempat parkir berisi $156$ mobil; $48$ pergi. Berapa yang tinggal?''
\end{exercise}

\begin{exercise}[$\star$]\label{exo:g2:problems:3}
Selesaikan: ``$5$ anak masing-masing memetik $4$ bunga. Berapa bunga
seluruhnya?''
\end{exercise}

\begin{exercise}[$\star$]\label{exo:g2:problems:4}
Gambarlah gambar batangnya, lalu selesaikan: ``Tono punya $45$ kartu,
kartu Mia $17$ lebih sedikit. Berapa kartu Mia?''
\end{exercise}

\begin{exercise}[$\star$]\label{exo:g2:problems:5}
Dua langkah: ``Satu peti berisi $6$ botol. Sebuah toko menerima $4$
peti lalu menjual $9$ botol. Berapa botol yang tersisa?''
\end{exercise}

\begin{exercise}[$\star$]\label{exo:g2:problems:6}
Carilah lalu coretlah bilangan yang tidak berguna, lalu selesaikan:
``Beni berumur $8$ tahun dan punya $23$ kelereng. Ia menang $15$
lagi. Berapa kelerengnya sekarang?''
\end{exercise}

\begin{exercise}[$\star$]\label{exo:g2:problems:7}
Dengan tabel pada \cref{ex:g2:problems:table}: berapa pai yang
terjual selama dua hari? Berapa kue dari segala jenis pada hari Sabtu?
\end{exercise}

\begin{exercise}[$\star$]\label{exo:g2:problems:8}
Dengan piktogram pada \cref{ex:g2:problems:pictogram}: berapa buah
seluruhnya? Warna mana yang paling banyak?
\end{exercise}

\begin{exercise}[$\star$]\label{exo:g2:problems:9}
Selesaikan: ``Sehelai pita panjangnya $80$ cm. Lia memotong $35$ cm.
Berapa panjang potongan yang tersisa?''
\end{exercise}

\begin{exercise}[$\star\star$]\label{exo:g2:problems:10}
``Seorang pembuat roti membuat $5$ loyang berisi $4$ kue pada pagi
hari dan $3$ loyang berisi $4$ kue pada sore hari. Berapa kue pada
hari itu?'' (Dua perkalian, lalu satu \omterm{def:g1:addition:def}{jumlah} --- atau bilanglah loyangnya dulu!)

\end{exercise}
